% AUTO-GENERATED by research/publication.py — sources: results/cyber_range/experiments/latest.json:suites.detector
\begin{table*}[t]
\centering
\footnotesize
\caption{Cyber Range: Which Detector Raises the Incident}
\label{tab:range_det}
\begin{tabular}{lccccc}
\toprule
\textbf{Scenario} & \textbf{Detector} & \textbf{Runs detected} & \textbf{Detection latency (s)} & \textbf{Attack-flow recall} & \textbf{Benign flow FPR} \\
\midrule
Zero-day intrusion & Frozen v1.1 & 0 / 10 & --- & 0.000 & 0.08\% \\
 & Zero-Day Gate & 10 / 10 & 4.3 [2.6, 6.0] & 0.645 & 0.81\% \\
Beacon-style bot & Frozen v1.1 & 10 / 10 & 4.1 [2.3, 5.9] & 0.563 & 0.09\% \\
 & Zero-Day Gate & 10 / 10 & 2.8 [1.2, 4.4] & 0.877 & 0.78\% \\
Web request burst & Frozen v1.1 & 10 / 10 & 2.6 [0.6, 4.6] & 0.907 & 0.09\% \\
 & Zero-Day Gate & 10 / 10 & 2.6 [0.6, 4.6] & 1.000 & 0.87\% \\
\bottomrule
\end{tabular}
\\[2pt]\parbox{0.95\linewidth}{\scriptsize 10 seeds per cell; incident rule 3 alerts within 5 s. Shown: scenarios where a detector missed runs or had attack-flow recall below 0.99. The other 6 scenarios were detected by both detectors in every run (mean latency 0.0--2.4 s). Mean false incidents per run across all cells: 0.01. The zero-day scenario redraws the same 47 never-seen flows, so its runs are not independent zero-days.}
\end{table*}
